\section{Related Work}
\label{sec:related}

\paragraph{Weight-only quantization.}
GPTQ uses approximate second-order information for accurate post-training
weight quantization~\citep{frantar2023gptq}.  AWQ instead identifies salient
weight channels from activation statistics and searches per-channel scaling,
providing strong low-bit accuracy and hardware-friendly inference
kernels~\citep{lin2024awq}.  Both are calibration-dependent in the sense relevant
to mixed-precision search: changing a layer-wise allocation can require an
expensive reconstruction or calibration pass.  HQQ uses half-quadratic splitting
to quantize weights without activation calibration~\citep{badri2023hqq}.  Its
calibration-free construction makes it possible to precompute a bank for each
bit-width and assemble arbitrary layer-wise configurations cheaply, which is why
we use HQQ as a search proxy rather than as the final deployment objective.

Mixed-precision weight methods allocate additional precision to sensitive
layers, channels, or outliers.  These methods can achieve better quality at a
fractional average bit-width, but fine-grained formats may introduce irregular
memory access or specialized-kernel requirements.  Our search uses one precision
per linear layer, following the hardware-compatible granularity emphasized by
AMQ~\citep{lee2025amq}.  The contribution here is not a new scalar weight
quantizer; it is the joint allocation of a deployment quantizer across weights
and a separately growing KV-cache budget.

\paragraph{Automated mixed-precision search.}
AMQ is the closest weight-side predecessor~\citep{lee2025amq}.  It prunes a
layer-wise search space with sensitivity, uses HQQ as a proxy for AWQ/GPTQ,
predicts Jensen--Shannon divergence with an RBF model, and iteratively updates an
NSGA-II archive.  AMQ also proves that global order equivalence between proxy and
deployment scores gives identical Pareto frontiers.  We retain the useful
per-axis search machinery but address a setting in which global order equivalence
does not hold in the joint interior.  Our theory therefore weakens the required
condition for \emph{coverage} to small conditional ordering violations, and our
algorithm reserves deployment-method measurements for the decisions that the
proxy cannot resolve.

Surrogate-assisted neural architecture search replaces many costly evaluations
with learned performance predictors~\citep{baker2017nas,white2021predictors}.
Multi-objective evolutionary algorithms such as NSGA-II and NSGA-III maintain a
set of accuracy--efficiency trade-offs rather than optimizing one fixed
budget~\citep{deb2002nsga2,deb2014nsga3}.  In contrast to a predictor trained
directly on hundreds of discrete layer variables, we use supervised dimension
reduction learned from the much larger proxy archives.  PLS is particularly
appropriate because it chooses latent directions according to covariance with
the target rather than input variance alone~\citep{wold2001pls}.  The deployment
labels still fit the final surrogate, preventing proxy values from being treated
as deployment ground truth.

\paragraph{KV-cache compression.}
KIVI observes different distributions for keys and values and applies
per-channel key quantization and per-token value quantization, together with a
full-precision residual window~\citep{liu2024kivi}.  KVQuant likewise studies
outliers, sensitivity-aware datatypes, and low-bit KV-cache
representations~\citep{hooper2024kvquant}.  ThinK compresses the channel dimension
by query-driven pruning and can be composed with KIVI~\citep{xu2025think}.  These
works supply the quantization and pruning primitives in our effective-KV axis.
They do not decide how a shared memory budget should be exchanged against
layer-wise weight precision.

KVTuner searches hardware-friendly layer-wise K/V precision pairs after
sensitivity-based pruning and clustering~\citep{li2025kvtuner}.  It is closely
related to our Stage~1 effective-KV search.  Our setting additionally couples
K/V bit-width and group size with retained channel dimension, and then combines
that axis with a layer-wise weight allocation.  This creates two new issues: the
Cartesian product of the axes is enormous, and the quality ordering within that
product depends on whether the weight quantizer is activation-independent or
activation-aware.  Our front-product theorem and proxy-boundary analysis target
these cross-axis issues.

\paragraph{Positioning.}
Prior work commonly uses sensitivity or Pareto pruning \emph{within} a
quantization axis.  Our result concerns a different intermediate regime: the
resource accounting is separable across weight and KV components, but the
quality objective is not assumed to be separable.  The product-front restriction
is justified by an audited conditional order, not by assuming that weight and KV
losses add.  This distinction is essential because the observed joint loss is
nearly additive in the deployable region yet exhibits substantial saturation in
the aggressive low-weight/low-KV corner.
